\documentclass{article}
\usepackage[utf8]{inputenc}
\usepackage{amsmath}
\usepackage{graphicx}
\usepackage{hyperref}
\usepackage{tcolorbox}
\usepackage[margin=1in]{geometry}
\usepackage{subcaption}
\usepackage{url}
\usepackage{float}
\usepackage{listings}
\usepackage{color}
\usepackage{tikz}

%%%%%%%%%%%%%%%%%%%%%%% HEADER %%%%%%%%%%%%%%%%%%%%%%%%

\title{Event Based Monte Carlo Using Pumi-PIC}
\author{Fuad Hasan}
\date{\today}

\begin{document}

\maketitle

%%%%%%%%%%%%%%%%%%%%% Mesh Description %%%%%%%%%%%%%%%%

\section{Mesh Description}
For the DAGMC simulations, two different meshes are used. DAGMC uses an \textit{unstructured surface mesh} for particle transport
and a \textit{volume mesh} for tallying flux using the \textit{track-length} estimator.
When comparing the performance of DAGMC with pumipush, two different volume meshes are used: a coarse and a fine mesh.
The mesh properties and the side lengths are shown in Table \ref{tab:mesh_properties}.

\begin{table}[h!]
	\centering
	\begin{tabular}{l|l}
		\hline
		\textbf{Parameter}                                 & \textbf{Value}                                  \\  \hline
		\hline
		Mesh Type                                          & Unstructured Tris and Tets                      \\
		\hline
		Height, Width, Depth                               & -0.5 to 0.5 (1) -0.5 to 0.5 (1) -0.5 to 0.5 (1) \\
		\hline
		\textit{Fine Mesh}: \# Triangles in Surface Mesh   & 1,404                                           \\
		\hline
		\textit{Fine Mesh}: \# Tetrahedra in Volume Mesh   & \textbf{97,410}                                 \\
		\hline
		\textit{Fine Mesh}: \# Nodes                       & 18,162                                          \\
		\hline
		\textit{Coarse Mesh}: \# Tetrahedra in Volume Mesh & \textbf{8,805}                                  \\
		\hline
		\textit{Coarse Mesh}: \# Nodes                     & 1,853                                           \\
		\hline
		\hline
	\end{tabular}
	\caption{Mesh Properties}
	\label{tab:mesh_properties}
\end{table}

The fine and coarse meshes are shown in Figure \ref{fig:meshes}.
\begin{figure}[h!]
	\centering
	\begin{subfigure}{0.50\textwidth}
		\centering
		\includegraphics[width=\textwidth]{images/finemesh3d.png}
		\caption{Fine Mesh}
		\label{fig:fine_mesh}
	\end{subfigure}
	\hfill
	\begin{subfigure}{0.45\textwidth}
		\centering
		\includegraphics[width=\textwidth]{images/coarsemesh3d.png}
		\caption{Coarse Mesh}
		\label{fig:coarse_mesh}
	\end{subfigure}
	\caption{Fine and Coarse Meshes}
	\label{fig:meshes}
\end{figure}

\subsubsection{Mesh Conversion Process}
Both the surface mesh and the volume mesh are created using Coreform Cubit. For the surface mesh, it is directly exported
to the \texttt{.h5m} format
and used in DAGMC neutron transport. This mesh needs material information and boundary conditions. These are set up following
the DAGMC documentation and this Coreform Cubit tutorial here:
\begin{tcolorbox}[colback=white, colframe=black, title=Cubit Tutorial]
	\url{https://youtu.be/2TzgTQidfwk}
\end{tcolorbox}

For the volume mesh, it can be created in any meshing software and DAGMC's \texttt{mbconvert} utility can be used to convert it to
both \texttt{.h5m} and \texttt{.msh} formats for DAGMC (\texttt{h5m}) and pumipush (\texttt{msh}), respectively. One caveat is that
the volume mesh seems to be
not working directly for pumipush (the internal gmsh reader in Pumi-PIC couldn't read it). I had to reload it in gmsh and save it
in ascii version 2 format to be read by pumipush. The following command is used to convert the volume mesh:
\begin{tcolorbox}[colback=white, colframe=black, title=Mesh Conversion Command]
	\begin{verbatim}
		gmsh -o cube8k.v2-gmsh.msh -save -3 -order 1 -save_all -format msh2 cube8k.v2.msh
	\end{verbatim}
\end{tcolorbox}

In this case, the volume mesh was created in Cubit and exported as \texttt{.e} (Exodus) format and then converted to \texttt{.msh} format
using \texttt{mbconvert} utility. The native \texttt{.cub} format worked before but with the new mesh, it was not working.

\section{Verification}
To verify that DAGMC and pumipush are simulating the same problem, we need to verify that equvalent computations are
being done in both codes. There are multiple facets to this verification:

\begin{itemize}
	\item \textbf{Mesh}: The above section shows that the meshes used in both codes are the same. Although, for transport, DAGMC uses
	      the surface mesh whereas pumipush uses the volume mesh for both transport and tallying. But this is not against the essence of this study.
	      Later, it is also shown that even if the surface mesh is replaced by the CSG geometry, the performance remains almost the same.
	\item \textbf{Material Properties}: In this simulation, a simple material with $\Sigma_t = 100.0$ and $\Sigma_s = 100.0$ and then
	      $\Sigma_t = 10.0$ and $\Sigma_s = 10.0$ are used. For pumipush, $\lambda = 0.01$ and $\lambda = 0.1$ are respectively for equivalency.
	      This means that only scattering is happening. All the neutrons are in one group. In pumipush, it used mean free path to define
	      the cross sections. Mean free path is defined as the inverse of the total cross section.
	\item \textbf{Source}: A fixed source simulation is done here. The source is uniformly distributed in the entire cube.
	      The source intensity can be different becuase we will compare the flux shape, not the nominal value. Moreover, the number of
	      particles being simulated is given as an input parameter for both of them.
	\item \textbf{Tally}: Only the flux tally is used in the performance analysis. Both of them use the "track-length estimator" (more
	      explanation later).
	\item \textbf{Event Based Method}: Paricles start their journey together and tracked until they either reach their destination or
	      get leaked out. For both cases, \textit{vacuum boundary} condition is used.
\end{itemize}

\subsection{Mean Free Path}
From the OpenMC documentation, the probability of a particle traveling a distance $\ell$ without interaction is given by:
\begin{equation}
	p(\ell) d \ell=\Sigma_t e^{-\Sigma_t \ell} d \ell
\end{equation}

where $\Sigma_t$ is the total cross section. Some exponential distributions are shown in Figure \ref{fig:exponential_distributions} as
a visual representation of the above equation.

\begin{figure}[h!]
	\centering
	\includegraphics[width=0.6\textwidth]{images/Exponential_distribution_pdf_-_public_domain.svg.png}
	\caption{Exponential Distributions}
	\label{fig:exponential_distributions}
\end{figure}

Therefore, the random distance $\ell$ is sampled using an uniform random number $\xi$
as follows:
\begin{equation}
	\ell=-\frac{\ln \xi}{\Sigma_t}
\end{equation}
$\xi$ is sampled from a uniform distribution between 0 and 1.
\\

In the pumipush code, it is sampled as follows:
\begin{verbatim}
    double rn = gen.drand(0.0, 1.0);
    double distance = -Kokkos::log(rn) * lambda;
\end{verbatim}
where \texttt{rn} is the random number sampled from uniform distribution from 0 to 1.
\\

In OpenMC, the group cross sections are specified as follows:
\begin{verbatim}
    groups = openmc.mgxs.EnergyGroups((0,1))
    xsdata = openmc.XSdata('steel', groups)
    xsdata.order = 0

    scatter_xs = 100.
    xsdata.set_total([scatter_xs])
    xsdata.set_absorption([0.])
    xsdata.set_scatter_matrix(np.array([scatter_xs]).reshape(1,1,1))

    mg_xs_file = openmc.MGXSLibrary(groups)
    mg_xs_file.add_xsdata(xsdata)
\end{verbatim}
This means: a single group cross section is used where the cross section remains constant at 100.0. $\Sigma_t = \Sigma_s = 100.0$, i.e.,
only scattering is happening and no absorption. Therefore, the mean free path ($\lambda$) is given by:
\begin{equation}
	\lambda = \frac{1}{\Sigma_t} = \frac{1}{100.0} = 0.01
\end{equation}

To confim that OpenMC is using this given total and scattering cross section, the results are analyzed. For verification, an
extra tally is added (it is removed during performance analysis) to tally the total flux, total reaction rate, and total scattering
reaction rate. There is a simple relation between flux and reaction rates:

\begin{equation}
	R_x = \Sigma_x \times \text{Flux}
\end{equation}

This implies that the $\frac{R_{\text{scattering}}}{\text{Flux}}$ will be equal to $\Sigma_s$ or 100.0. It will be the same for
$\frac{R_{\text{total}}}{\text{Flux}}$ and $\Sigma_t$. For absorption, it will be zero.
\\

These quantities are found as follows in OpenMC:
\begin{align}
	\frac{R_{\text{scattering}}}{\text{Flux}} & = 100.00000000000517                 \\
	\frac{R_{\text{total}}}{\text{Flux}}      & = 100.00000000000517                 \\
	\frac{R_{\text{absorption}}}{\text{Flux}} & = 1.0000000000000523 \times 10^{-10}
\end{align}
In the above simulation 20,000 particles are used. The small difference is due to the statistical error.



\subsection{Event Based Method}
Event based method in OpenMC is implemented here:
\begin{tcolorbox}[colback=white, colframe=black, title=Event Based Method Code]
	\url{https://github.com/openmc-dev/openmc/blob/d30b2e801413faaee07282969955991abdf7baeb/src/simulation.cpp#L799-L849}
\end{tcolorbox}

\subsubsection*{How it Works}
Roughly, the simulation works as follows:
\begin{enumerate}
	\item Determine/initialize which particle is going to do what if not already in flight in the current \textit{subiteration}.
	\item There are specific type of events and the particles are sorted based on the even type.
	\item All types of events happen until there is no event left.
	\item Remove particles that have died due to various reasons.
	\item Repeat the process until the number of particles are zero.
\end{enumerate}

% page break
\newpage
\section{Validation of PumiPush Transport and Tallying}
To properly validate that the simulation are producing correct results, flux is tallied in both codes. They are normalized
using the sum of the flux values and plotted in Figure \ref{fig:flux_comparison}.

\begin{figure}[h!]
	\centering
	\begin{subfigure}{0.45\textwidth}
		\centering
		\includegraphics[width=\textwidth]{images/dagmc_flux_yz_plane.png}
		\caption{DAGMC Computed Flux}
		\label{fig:flux_dagmc}
	\end{subfigure}
	\hfill
	\begin{subfigure}{0.45\textwidth}
		\centering
		\includegraphics[width=\textwidth]{images/pumipush_flux_yz_plane.png}
		\caption{PumiPush Computed Flux}
		\label{fig:flux_pumipush}
	\end{subfigure}
	\caption{Flux on $x=0$ Plane}
	\label{fig:flux_comparison}
\end{figure}

\begin{figure}[h!]
	\centering
	\includegraphics[width=0.8\textwidth]{images/flux_along_y_line.png}
	\caption{Flux comparison along z-axis}
	\label{fig:flux_comparison_along_y_axis}
\end{figure}

The flux along a line that passes through the center of the cube along the y axis is plotted in Figure \ref{fig:flux_comparison_along_y_axis}.
This is drawn with the fine mesh with $\lambda = 0.01$ and 2,000,000 particles.
It shows that the flux distributions are similar in both codes. This makes sure that the
pumipic track-length estimator implementation is correct.


\section{Used Tools}
\begin{table}[h!]
	\centering
	\begin{tabular}{l|l}
		\hline
		\textbf{Tool} & \textbf{Version}        \\ \hline \hline
		Pumi-PIC      & commit \texttt{e1b04da} \\ \hline
		OpenMC        & 0.15.0 through mamba    \\ \hline
		GPU           & A100 on Perlmutter      \\ \hline
	\end{tabular}
	\caption{Software Versions}
	\label{tab:software_versions}
\end{table}


\newpage
\section{Performance Analysis}
The amount of work for the simulation to be completed is dependent on the following factors: (i) the number of particles being simulated,
(ii) the average distance each particle travels (since the higher the distance, the more mesh elements it will cross), and (iii) the mesh:
coarseness and total number of elements.


\subsection{Fine Mesh}
In the following Figure \ref{fig:performance_analysis01}, the time taken to simulate for $\Sigma_t = 100.0$ or $\lambda = 0.01$ is shown
for varying number of particles.
\begin{figure}[h!]
	\centering
	\begin{subfigure}{0.49\textwidth}
		\centering
		\includegraphics[width=\textwidth]{images/pumipic_vs_dagmc0.01.png}
		\caption{Performance Analysis for $\lambda = 0.01$}
		\label{fig:performance_analysis01}
	\end{subfigure}
	\hfill
	\begin{subfigure}{0.49\textwidth}
		\centering
		\includegraphics[width=\textwidth]{images/pumipic_vs_dagmc0.1.png}
		\caption{Performance Analysis for $\lambda = 0.1$}
		\label{fig:performance_analysis1}
	\end{subfigure}
	\caption{Performance Analysis on the Fine Mesh}
	\label{fig:performance_analysis_combined}
\end{figure}

For the case of $\lambda = 0.1$ (Figure \ref{fig:performance_analysis1}) the pumipush code is much faster than $0.01$ case. It is due to following reasons:
\begin{itemize}
	\item The particles are moving further in the mesh and the leakage rate is higher. Therefore, it needs less
	      number of subiterations to complete the simulation.
	\item For the same reason, the \texttt{search\_mesh} method has more work at once on GPU (to search the intersected
	      mesh faces and elements) and it is more efficient rather that multiple kernel launches.
\end{itemize}

The \texttt{DAGMC with CSG} line in Figure \ref{fig:performance_analysis1} shows the simulation time for the case where only the flux tally is computed on
the unstructured mesh and the transport is done using the CSG gemoetry rather than the surface mesh used before.
It implies that finding out where the particles are in the mesh is not the bottleneck for DAGMC but the mesh traversal for tallying
is.

\subsection{Coarse Mesh}
To understand how the scalability seen in the fine mesh case translates to the coarse mesh, the performance analysis is done on the
coarse mesh which is about 10 times coarser. Since, it is already seen for DAGMC that the performance is not affected significantly
when CSG geometry is used, I continued to used CSG geometry for the coarse mesh case.
The results are shown in Figure \ref{fig:performance_analysis_coarse}.

\begin{figure}[h!]
	\centering
	\begin{subfigure}{0.49\textwidth}
		\centering
		\includegraphics[width=\textwidth]{images/pumipic_vs_dagmc_coarse0.01.png}
		\caption{Simulation Time for $\lambda = 0.01$}
		\label{fig:performance_analysis01_coarse}
	\end{subfigure}
	\hfill
	\begin{subfigure}{0.49\textwidth}
		\centering
		\includegraphics[width=\textwidth]{images/pumipic_vs_dagmc_coarse0.1.png}
		\caption{Simulation Time for $\lambda = 0.1$}
		\label{fig:performance_analysis1_coarse}
	\end{subfigure}
	\caption{Performance Analysis on the Coarse Mesh}
	\label{fig:performance_analysis_coarse}
\end{figure}

The trends remain same for the coarse mesh as well. The speedup due to the coarseness of the mesh is shown in Figure \ref{fig:speedup}.
It is interesting that the speedup is not close to 10 but it is higher for the case of $\lambda = 0.1$ as expected due to higher number of
line face intersections.

\begin{figure}[h!]
	\centering
	\includegraphics[width=0.6\textwidth]{images/speedup_coarse_to_fine.png}
	\caption{Speedup Due to Coarseness of the Mesh}
	\label{fig:speedup}
\end{figure}

\newpage
\section{Future Works}
\begin{itemize}
	\item \textbf{Single Energy Group in DAGMC}: Now DAGMC is simulation with a single energy group but it has a range.
	      Although my understanding is that the particles shouldn't change energy during scattering here but it is worth checking.
	\item \textbf{Number of OpenMP Threads}: How to make sure that pumipush is using 128 threads? I have used the following command:
	      \begin{tcolorbox}[colback=white, colframe=green, title=OMP Environment Variables]
		      \begin{verbatim}
export OMP_PROC_BIND=spread && export OMP_PLACES=threads && 
export OMP_NUM_THREADS=128
    \end{verbatim}
	      \end{tcolorbox}
	\item \textbf{Memory Warning During DAGMC Run}: What does the following error mean?
	      \begin{tcolorbox}[colback=white, colframe=yellow, title=Memory Warning]
		      \small
		      \begin{verbatim}
[nid004261:462725] shmem: mmap: an error occurred while determining whether or not 
/tmp/ompi.nid004261.104431/jf.0/3845718016/shared_mem_cuda_pool.nid004261 could be created.
[nid004261:462725] create_and_attach: unable to create shared memory 
BTL coordinating structure :: size 134217728
        \end{verbatim}
	      \end{tcolorbox}
\end{itemize}

\bibliographystyle{plain}
\bibliography{references}

\end{document}